\documentclass[11pt]{article}
\usepackage{amsmath,amssymb,amsthm,hyperref}
\newtheorem{proposition}{Proposition}
\newtheorem{lemma}{Lemma}
\title{From overlapping local estimates to actual rectangular means:
energy transfer and the common-refinement count obstruction}
\author{The Eriksson Programme --- grouped research dossier}
\date{5 November 2026; unit 132; unrendered draft source}
\begin{document}
\maketitle
\paragraph{Purpose and status.}
We separate two obligations in the single coefficient-adapted rectangular
observation question: transferring energy estimates and keeping a linear
cell count. The first has an elementary answer below. The second does not
follow from bounded overlap through a common refinement. These are finite
analytic statements, with proofs and preparing-agent self-audit only. No
formal check, executed diagnostic, external review, originality claim,
new paper, or physical/continuum conclusion is asserted.

\paragraph{Conventions.}
Let $T_n=(\mathbb Z/n\mathbb Z)^D$, $n\ge3$, with counting measure and
unit-spacing forward edges $(x,i)$ from $x$ to $x+e_i$. The coefficient
$a(x)>0$ is shared among directions. Fields take values in a finite-dimensional
real Hilbert space $H$. Define
\[
 E_a(u)=\sum_{x,i}a(x)\|u(x+e_i)-u(x)\|^2,
 \quad V_B(u)=\sum_{x\in B}\|u(x)-\bar u_B\|^2,
\]
where $\bar u_B=|B|^{-1}\sum_Bu$. The internal energy $E_{a,B}$ includes
precisely those forward edge occurrences whose two endpoints belong to $B$.
For full cyclic coordinates it includes their wrap edges. Every cell in a
rectangular partition is a product of nonempty cyclic intervals. The actual
normalized mean map is $C_{\cal P}u=(|Q|^{-1/2}\sum_Qu)_{Q\in\cal P}$.

\begin{proposition}[Transfer to any subordinate partition]
Suppose nonempty sets $B_1,\ldots,B_m$ cover $T_n$, every vertex belongs to
at most $r$ of them, and for all fields
\[
 \kappa V_{B_j}(u)\le E_{a,B_j}(u),\qquad j=1,\ldots,m.
\]
If a partition $\cal P$ admits an assignment $j(Q)$ with $Q\subset B_{j(Q)}$,
then
\[
 \sum_{Q\in\cal P} V_Q(u)\le\frac r\kappa E_a(u).
\]
Consequently, if $\kappa\ge rc$, its actual means satisfy
\[
 c\|u\|_2^2\le E_a(u)+c\|C_{\cal P}u\|_2^2.
\]
The cells need not have separately proved internal Poincare inequalities.
The conclusion holds for the same fixed coefficient as in the local bounds.
\end{proposition}
\begin{proof}
The arithmetic mean minimizes squared distance to constants, so
\[
 \sum_{Q:j(Q)=j}V_Q(u)
 \le\sum_{Q:j(Q)=j}\sum_{x\in Q}\|u(x)-\bar u_{B_j}\|^2
 \le V_{B_j}(u).
\]
Assigned cells are disjoint, and all the omitted summands are nonnegative.
Summing the assumed local estimates gives
\[
 \kappa\sum_QV_Q(u)\le\sum_jE_{a,B_j}(u)\le rE_a(u).
\]
Indeed any directed edge included in $E_{a,B_j}$ has its tail in $B_j$,
so it is counted at most $r$ times. Nonnegativity is essential at this step.
Finally orthogonal block-mean decomposition gives exactly
$\|u\|_2^2=\sum_QV_Q(u)+\|C_{\cal P}u\|_2^2$.
\end{proof}

\paragraph{What was and was not eliminated.}
The proof derives the transfer; it does not assume a new transfer structure.
It also shows why discarding all cross-cell energies and demanding an
independent local estimate on every final cell is unnecessarily strong.
Cover-internal edges joining different final cells are retained in the
intermediate sum. However the existence, count and local constants of an
appropriate cover have not been derived from the original critical moment.
Those are still hypotheses of this conditional lemma, not hypotheses removed
from the original theorem. A merely scalar local estimate can be extended
componentwise to $H$ by summing in an orthonormal basis; no assertion about an
interacting gauge energy follows.

\begin{proposition}[Unconditional geometric conversion, with a power loss]
If the cover above consists of $m$ cyclic axis-parallel rectangles, there is
an actual rectangular partition subordinate to the cover with
\[
 \#\cal P\le \min\{n^D,(2m+1)^D\}.
\]
It can be chosen to refine every cover member. The transfer proposition
applies, but no linear bound in $m$ follows from this construction.
\end{proposition}
\begin{proof}
Represent each cyclic coordinate as $\{0,\ldots,n-1\}$. For each direction,
cut at $0,n$ and at the at most two endpoints of each proper cyclic
interval occurring in a cover member; a wrapped interval is the union of the
two corresponding linear end intervals. Whole cycles require no endpoints.
There are at most $2m+1$ nonempty coordinate segments. Their Cartesian
products form a partition of at most $(2m+1)^D$ nonempty rectangles, and at
most $n^D$ because every cell contains a vertex. Each cover member is a union
of grid cells. Since it is a cover, every grid cell belongs to at least one
cover member; assign one of them by a fixed index rule. The rule is geometric
and independent of the field.
\end{proof}

\paragraph{Count relevance.}
If one had a bounded-overlap cover with
$m\le A\max\{1,n^D/(1+\mu)^{D/2}\}$ and local constant $\kappa\ge rc$,
the power-loss conversion proves a corresponding power-law cell count.
It does \emph{not} prove the target linear order when that target grows.
When the target stays bounded, a fixed power of its bound is still bounded;
the count objection alone does not rule out that regime. Neither the
existence of this hypothetical cover nor a linear subordinate conversion is
proved here.

\begin{proposition}[Depth-three common refinements can be quadratic]
For every integer $J\ge2$ and every $D\ge2$, let $n=2J$. There is a cover
of $T_n$ by $m=2J+1$ axis-parallel rectangles with vertex overlap at most
three such that every rectangular partition refining \emph{all} cover
members has at least
\[
 n^2=4J^2=(m-1)^2
\]
cells. This lower bound is attained. In contrast, a partition merely
subordinate to the same cover can have one cell.
\end{proposition}
\begin{proof}
Take the full torus $B_0=T_n$, the $J$ strips
$V_j=\{x:x_1=2j+1\}$, and the $J$ strips
$H_j=\{x:x_2=2j+1\}$, for $j=0,\ldots,J-1$.
Each vertex lies in $B_0$, at most one $V_j$ and at most one $H_j$.

Refining a set means every partition cell is wholly inside it or disjoint
from it. A nonempty cell whose first cyclic interval contains at least two
vertices contains a consecutive odd/even pair (also at the cyclic seam,
because $n$ is even). It then intersects the corresponding $V_j$ but is
not contained in it. This contradicts refinement. Thus its first interval
is a singleton. The same argument using the $H_j$ makes its second interval
a singleton. Every cell has at most $n^{D-2}$ vertices, giving at least
$n^D/n^{D-2}=n^2$ cells. All products
$\{x_1\}\times\{x_2\}\times(\mathbb Z/n\mathbb Z)^{D-2}$ attain the bound.
For mere subordination, the one-cell partition $\{T_n\}$ is assigned to
$B_0$ and is valid. Hence the example does not disprove the existence of
some linear-size subordinate partition; it only disproves a linear-size
universal common refinement.
\end{proof}

\paragraph{A local-energy control for the geometric example.}
This example need not rely on a cover with invalid local estimates. Fix
$c>0$ and take a constant coefficient
$a(x)=\mu=\max\{1,3cn^2/16\}$. Every member is a product of full cycles
and singleton coordinates, with at least one full cycle. Its nonconstant
Fourier frequencies have energy eigenvalue at least
$4\mu\sin^2(\pi/n)\ge16\mu/n^2\ge3c$.
Thus all cover-local estimates hold with $\kappa=3c$ and $r=3$.
The original coefficient range, exact mean and critical moment hold, the
last with $M=1$. Yet common refinement still wastes quadratically many
cells in the cover size, whereas the full-torus member already supplies a
one-box estimate. This demonstrates waste in the conversion method, not a
coefficient counterexample to the original question; this particular cover
is not alleged to have optimal size relative to the moment target.

\paragraph{Primary-source comparison and originality.}
Sukochev--Zanin~\cite{SZ22} studies weak Schatten estimates for a
Birman--Schwinger operator with smoothing order $D/2$ and an $L\log L$
multiplier. This is useful context for critical spectral methods, but its
operator and hypotheses differ from the present weighted first-difference
rectangular-means target. No theorem from that paper is used as a local
cover certificate or a disjoint-partition theorem. The finite arguments
above are proved in full. Bounded-overlap energy counting and rectangular
arrangements are elementary; priority is unestablished and no independent
new theorem of the literature is claimed. Attempts to retrieve later
sections of the primary web text timed out; only its abstract/introduction
operator statement is relied on here.

\paragraph{Model assessments and review context.}
Preparing-agent self-audit only, in the same Codex session that derived these
claims. The system describes a GPT-6-based agent, while the actual served
backend identifier and reasoning setting are unexposed. No independent AI,
external service, human peer review, grade, PDF rendering, Lean run or
executed diagnostic is claimed. Checks derived the variance minimization,
edge-occurrence multiplicity, cyclic cuts, even-seam parity, attained cell
count, conditional constants and scope distinctions. No result or assessment
is inherited from another artifact.

\paragraph{Route closure and next unit.}
This closes a precise shortcut negatively: replacing a bounded-overlap
anisotropic cover by a common refinement cannot generally preserve linear
count, even at depth three and with valid local energy estimates. Do not
repeat this mechanism or relabel its power-loss count as the target theorem.
Simple subordination, special covers, or a coefficient-specific global
argument remain possible and are not refuted. The original fully variable
coupled divergence-zero question remains open.

Units 130--132 have not advanced that original question: finite verification,
classical quantifier audit, and now the explicit conversion boundary. Count
three conservatively; stop the generic cover/common-refinement research
mechanism and preserve this concrete blockage. The next work unit switches
to consolidating the already derived invariant-coordinate positive theorem
(unit 127, with actual finite diagnostics in unit 128), together with the
related layered/separated benchmarks and carefully scoped negative results.
Audit the exact combined manuscript and its primary-source novelty before
rendering/review/publication. Do not split these guards into micro-papers,
redeposit published v5, inherit grades, or open a second active mathematical
question. No original hypothesis is counted as eliminated.

\begin{thebibliography}{1}
\bibitem{SZ22} F. A. Sukochev and D. V. Zanin,
\emph{Solomyak-type eigenvalue estimates for the Birman--Schwinger operator},
Sb. Math. 213 (2022), 1250--1289.
\href{https://doi.org/10.4213/sm9732e}{doi:10.4213/sm9732e}.
\href{https://www.mathnet.ru/php/archive.phtml?jrnid=sm\&option_lang=eng\&paperid=9732\&wshow=paper}{Primary text}.
\end{thebibliography}
\end{document}
